\chapter{黑板推导训练：十条必须能重建的主线}

\begin{mapbox}
本章不是新的知识清单，而是闭卷训练脚本。每条推导都按“起点、关键步骤、检查、追问”组织。建议给自己 8 分钟，只写决定性步骤；写完再与正文核对。
\end{mapbox}

\section{Bloch 定理}
\textbf{起点：}$V(\vb r+\vb R)=V(\vb r)$，所以
$[H,T_{\vb R}]=0$。

\textbf{关键步骤：}
\[
T_{\vb R}|\psi\rangle=\lambda(\vb R)|\psi\rangle,\quad
|\lambda|=1,
\]
\[
\lambda(\vb R+\vb R')=\lambda(\vb R)\lambda(\vb R')
\Rightarrow\lambda=e^{i\vb k\cdot\vb R}.
\]
定义 $u_{\vb k}=e^{-i\vb k\cdot\vb r}\psi_{\vb k}$，验证
$u(\vb r+\vb R)=u(\vb r)$。

\textbf{检查：}自由电子时 $u$ 可取常数；$\vb k\to\vb k+\vb G$ 给同一平移本征值。

\textbf{追问：}准动量与机械动量有何区别？磁场下普通平移为何需要修改？

\section{近自由电子区边界开隙}
\textbf{起点：}$E_{\vb k}^{(0)}=E_{\vb k-\vb G}^{(0)}$ 在 Bragg 面成立。

\textbf{关键步骤：}
\[
\begin{pmatrix}
E_0+V_0&V_{\vb G}\\
V_{\vb G}^*&E_0+V_0
\end{pmatrix}
\binom{c_1}{c_2}=E\binom{c_1}{c_2},
\]
\[
E_\pm=E_0+V_0\pm|V_{\vb G}|,\qquad
\Delta E=2|V_{\vb G}|.
\]
\textbf{检查：}$V_G\to0$ 恢复交叉；若对称性使 $V_G=0$，可以不打开隙。

\textbf{追问：}驻波图像是什么？二维/三维中哪些面满足简并？

\section{Debye 热容}
\textbf{起点：}三条声学支，长波 $\omega=v_sq$。
\[
N(q)=3\frac{V}{(2\pi)^3}\frac{4\pi q^3}{3},
\qquad g(\omega)\propto\omega^2.
\]
\[
U=\int_0^{\omega_D}\frac{\hbar\omega\,g(\omega)}
{e^{\beta\hbar\omega}-1}\dd\omega.
\]
低温令 $x=\hbar\omega/\kb T$，得到 $U\propto T^4$、
$C_V\propto T^3$。

\textbf{检查：}高温回到 $3N\kb$；二维线性色散给 $T^2$。

\textbf{追问：}为什么 Einstein 模型低温失败？光学支低温贡献如何？

\section{Sommerfeld 电子热容}
\textbf{起点：}
$U=\int Eg(E)f(E)\dd E$，固定 $N$。

\textbf{关键步骤：}
\[
\int\phi(E)f(E)\dd E
=\int_0^\mu\phi(E)\dd E
+\frac{\pi^2}{6}(\kb T)^2\phi'(\mu)+\cdots.
\]
同时展开粒子数以求 $\mu(T)$，最终
\[
U(T)-U(0)=\frac{\pi^2}{6}g(E_F)(\kb T)^2,\quad
C_V=\frac{\pi^2}{3}\kb^2g(E_F)T.
\]
\textbf{检查：}经典均分给 $O(N\kb)$，量子结果被
$T/T_F$ 抑制。

\textbf{追问：}相互作用如何通过有效质量改变 $\gamma$？

\section{Boltzmann 电导}
\textbf{起点：}均匀稳态
\[
-\frac e\hbar\vb E\cdot\nabla_{\vb k}f
=-\frac{f-f_0}{\tau}.
\]
\textbf{关键步骤：}
\[
\nabla_{\vb k}f_0=\hbar\vb v\,\partial_Ef_0,\quad
\delta f=e\tau\vb E\cdot\vb v(-\partial_Ef_0).
\]
\[
\sigma_{\alpha\beta}=e^2\int\frac{\dd^3k}{(2\pi)^3}
\tau v_\alpha v_\beta(-\partial_Ef_0).
\]
\textbf{检查：}抛物线单带、常数 $\tau$ 恢复
$ne^2\tau/m$。

\textbf{追问：}为什么小角散射对量子寿命和输运寿命权重不同？

\section{正则分布与涨落}
\textbf{起点：}系统与大热库总能固定，
\[
\ln\Omega_R(E_{\rm tot}-E_i)
\simeq\ln\Omega_R(E_{\rm tot})-\beta E_i.
\]
\[
p_i=e^{-\beta E_i}/Z.
\]
\textbf{关键步骤：}
\[
U=-\partial_\beta\ln Z,\quad
\partial_\beta^2\ln Z=\langle(\Delta E)^2\rangle,
\]
\[
\langle(\Delta E)^2\rangle=\kb T^2C_V.
\]
\textbf{检查：}热容正时涨落非负；相对涨落随 $N^{-1/2}$ 消失。

\textbf{追问：}负热容为何可在某些长程相互作用微正则系统出现？

\section{Hubbard 到 Heisenberg}
\textbf{起点：}半填充、$U\gg t$，低能子空间每点一个电子。

\textbf{关键步骤：}一阶跃迁进入能量高 $U$ 的双占据子空间；二阶有效 Hamiltonian
\[
H_{\rm eff}=-P H_tQ\frac1UQH_tP.
\]
比较自旋单态与三重态，得到
\[
H_{\rm eff}=J\sum_{\langle ij\rangle}
\left(\vb S_i\cdot\vb S_j-\frac14n_in_j\right),
\quad J=4t^2/U.
\]
\textbf{检查：}$t\to0$ 交换消失；$U$ 越大虚跃迁越弱。

\textbf{追问：}掺杂后为什么得到 $t$--$J$ 模型？多轨道 Hund 耦合怎样改变交换？

\section{Berry 曲率与 Chern 数}
\textbf{起点：}绝热占据带 $|u_n(\vb k)\rangle$。
\[
\mathcal A_n=i\langle u_n|\nabla_k u_n\rangle,\quad
\Omega_n=\nabla_k\times\mathcal A_n.
\]
规范变换使 $\mathcal A\to\mathcal A-\nabla\chi$，但
$\Omega$ 不变。
\[
C_n=\frac1{2\pi}\int_{\rm BZ}\Omega_{n,z}\dd^2k\in\mathbb Z.
\]
\textbf{检查：}时间反演使
$\Omega_n(\vb k)=-\Omega_{\bar n}(-\vb k)$，完整 Kramers 占据子空间总 Chern 数为零。

\textbf{追问：}为什么量子自旋 Hall 仍可非平庸？非 Abelian Berry 连接何时需要？

\section{BCS Bogoliubov 对角化}
\textbf{起点：}定义
$\Delta=-V\sum_k\langle c_{-k\downarrow}c_{k\uparrow}\rangle$。
\[
\mathcal H_k=
\begin{pmatrix}\xi_k&\Delta\\\Delta^*&-\xi_k\end{pmatrix}.
\]
\textbf{关键步骤：}
\[
\det(\mathcal H_k-E)=E^2-\xi_k^2-|\Delta|^2=0,
\]
\[
E_k=\sqrt{\xi_k^2+|\Delta|^2}.
\]
由异常平均值回代得到能隙方程；弱耦合积分有对数，故
$\Delta\sim\omega_D e^{-1/N(0)V}$。

\textbf{检查：}$\Delta\to0$ 恢复粒子/空穴谱；费米面处最小隙为 $|\Delta|$。

\textbf{追问：}节点配对时 $E_k=0$ 在哪里？规范相位如何联系超流？

\section{Maxwell 到光学响应}
\textbf{起点：}频域 Maxwell 与
$\vb J=\sigma(\omega)\vb E$。
\[
\nabla\times\nabla\times\vb E
=\frac{\omega^2}{c^2}\vb E+i\mu_0\omega\vb J.
\]
均匀横波得到
\[
k^2=\varepsilon(\omega)\frac{\omega^2}{c^2},
\qquad
\varepsilon(\omega)=1+\frac{i\sigma(\omega)}
{\varepsilon_0\omega}.
\]
代 Drude 电导，在无阻尼极限得
$\varepsilon=1-\omega_p^2/\omega^2$。

\textbf{检查：}$\omega\gg\omega_p$ 时趋近真空传播；
$\omega<\omega_p$ 时金属反射。

\textbf{追问：}带间跃迁如何加入？介电函数零点和损耗函数峰有什么关系？

\section{黑板自评量表}
每条推导按下列标准各记 0--2 分：
\begin{enumerate}
\item 是否先声明模型与适用条件；
\item 是否写出起点而非从结论倒推；
\item 是否保留决定物理的中间步骤；
\item 是否做量纲或特殊极限检查；
\item 是否能接住一个实验或失效条件追问。
\end{enumerate}
8 分以上才算“面试可用”；5--7 分说明会做题但表达链不完整；4 分以下应回到对应正文，而不是继续刷更多题。

